% ============================================================
%  Relatório IAM TP — MEI, Gestão de Identidade — ESTG-IPP
%  2025/2026 — Avelino Almeida (8220800@estg.ipp.pt)
% ============================================================
\documentclass[12pt,a4paper]{article}

% --- Encoding & Language ---
\usepackage[utf8]{inputenc}
\usepackage[T1]{fontenc}
\usepackage[portuguese]{babel}

% --- Layout ---
\usepackage[top=2.5cm, bottom=2.5cm, left=2.8cm, right=2.8cm]{geometry}
\usepackage{setspace}
\onehalfspacing

% --- Fonts ---
\usepackage{lmodern}

% --- Graphics & Figures ---
\usepackage{graphicx}
\usepackage{float}
\usepackage{caption}
\usepackage{subcaption}
\graphicspath{{../imagens/}{../../app/static/report-assets/}}

% --- Tables ---
\usepackage{booktabs}
\usepackage{tabularx}
\usepackage{array}
\usepackage{multirow}
\usepackage{longtable}
\usepackage[table,xcdraw]{xcolor}

% --- Code Listings ---
\usepackage{listings}
\usepackage{lstautogobble}
\definecolor{codebg}{RGB}{245,247,250}
\definecolor{codeframe}{RGB}{209,219,239}
\definecolor{codekw}{RGB}{37,99,235}
\definecolor{codecomment}{RGB}{107,114,128}
\definecolor{codestring}{RGB}{5,150,105}
\lstset{
  backgroundcolor=\color{codebg},
  frame=single,
  rulecolor=\color{codeframe},
  basicstyle=\ttfamily\small,
  keywordstyle=\color{codekw}\bfseries,
  commentstyle=\color{codecomment}\itshape,
  stringstyle=\color{codestring},
  breaklines=true,
  breakatwhitespace=true,
  numbers=left,
  numberstyle=\tiny\color{gray},
  numbersep=8pt,
  showstringspaces=false,
  autogobble=true,
  xleftmargin=12pt,
  xrightmargin=4pt,
  aboveskip=10pt,
  belowskip=10pt,
}

% --- Hyperlinks ---
\usepackage[hidelinks,colorlinks=false]{hyperref}
\usepackage{url}

% --- Header/Footer ---
\usepackage{fancyhdr}
\pagestyle{fancy}
\fancyhf{}
\fancyhead[L]{\small\textit{IAM TP — Gestão de Identidade}}
\fancyhead[R]{\small\textit{ESTG-IPP · MEI · 2025/2026}}
\fancyfoot[C]{\thepage}
\renewcommand{\headrulewidth}{0.4pt}

% --- Section styling ---
\usepackage{titlesec}
\titleformat{\section}{\Large\bfseries\color{black}}{}{0em}{}[\titlerule]
\titleformat{\subsection}{\large\bfseries}{}{0em}{}
\titleformat{\subsubsection}{\normalsize\bfseries}{}{0em}{}

% --- Misc ---
\usepackage{enumitem}
\usepackage{mdframed}
\usepackage{tcolorbox}
\tcbuselibrary{skins,breakable}

% --- Callout box ---
\newtcolorbox{callout}[1][]{
  colback=blue!5!white,
  colframe=blue!60!white,
  fonttitle=\bfseries,
  title=#1,
  breakable,
  left=6pt, right=6pt, top=4pt, bottom=4pt
}
\newtcolorbox{warningbox}[1][]{
  colback=orange!8!white,
  colframe=orange!70!white,
  fonttitle=\bfseries,
  title=#1,
  breakable,
  left=6pt, right=6pt, top=4pt, bottom=4pt
}

% ============================================================
\begin{document}

% ============================================================
%  CAPA
% ============================================================
\begin{titlepage}
  \thispagestyle{empty}
  \centering
  \vspace*{1cm}

  {\large \textbf{Instituto Politécnico do Porto}}\\[0.3cm]
  {\large Escola Superior de Tecnologia e Gestão}\\[0.3cm]
  {\large Mestrado em Engenharia Informática}\\[2cm]

  \rule{\linewidth}{1.5pt}\\[0.5cm]
  {\Huge \bfseries Projeto de Gestão de Identidade}\\[0.4cm]
  {\LARGE \bfseries Plataforma IAM com Keycloak, FastAPI e Docker}\\[0.3cm]
  \rule{\linewidth}{1.5pt}\\[1.5cm]

  {\large \textbf{Unidade Curricular:} Gestão de Identidade}\\[0.3cm]
  {\large \textbf{Docente:} Ricardo Costa}\\[1.5cm]

  {\large \textbf{Autor:}}\\[0.3cm]
  {\Large Avelino Almeida — 8220800}\\[0.2cm]
  {\large 8220800@estg.ipp.pt}\\[2cm]

  {\large Ano letivo: 2025/2026}\\[0.5cm]
  {\large Data de entrega: 28 de maio de 2026}\\[2cm]

  \vfill
  {\small Relatório do Trabalho Prático — MEI-GI}
\end{titlepage}

% ============================================================
%  RESUMO
% ============================================================
\newpage
\thispagestyle{empty}
\section*{Resumo}

Este relatório descreve a concepção e implementação de uma plataforma de \textit{Identity and Access Management} (IAM) para o trabalho prático da unidade curricular de Gestão de Identidade do Mestrado em Engenharia Informática (ESTG-IPP, 2025/2026).

A solução é totalmente containerizada e combina o \textbf{Keycloak 26} como \textit{Identity Provider} (IdP) central, uma \textbf{API FastAPI} (Python 3.12) como servidor de recursos protegido e \textbf{PostgreSQL 16} como base de dados de suporte. O projeto demonstra os seguintes pilares de IAM:

\begin{itemize}[leftmargin=1.5cm]
  \item \textbf{Autenticação federada} via OIDC/OAuth 2.0, com validação de JWT RS256 por JWKS;
  \item \textbf{Autorização RBAC} com três papéis hierárquicos (\textit{admin}, \textit{colaborador}, \textit{visitante});
  \item \textbf{Autenticação Multifator (MFA/TOTP)} obrigatória para administradores;
  \item \textbf{Ciclo de vida de identidades (JML)}: Joiner, Mover e Leaver via Admin REST API;
  \item \textbf{Auditoria} de eventos de autenticação e administração.
\end{itemize}

A infraestrutura arranca com um único comando (\texttt{docker compose up --build}) e inclui um dashboard interativo para demonstração de todos os requisitos.

\vspace{0.5cm}
\noindent\textbf{Palavras-chave:} IAM, Identity Management, Keycloak, OIDC, OAuth 2.0, JWT, RBAC, MFA, TOTP, JML, FastAPI, Docker.

\newpage
\tableofcontents

% ============================================================
%  1. INTRODUÇÃO
% ============================================================
\newpage
\section{Introdução e Objetivos}

\subsection{Contexto e Motivação}

A gestão de identidades e acessos (\textit{Identity and Access Management} — IAM) constitui um dos pilares fundamentais da segurança empresarial moderna. À medida que as organizações adoptam arquitecturas distribuídas, \textit{cloud} e modelos de trabalho híbrido, a necessidade de um sistema centralizado que gerencie \textit{quem} pode aceder a \textit{quê}, \textit{quando} e \textit{como} torna-se crítica.

Este trabalho prático tem como objectivo construir uma solução IAM funcional, reproduzível e auditável, que demonstre a aplicação prática dos conceitos teóricos da unidade curricular de Gestão de Identidade.

\subsection{Objetivos Específicos}

\begin{enumerate}[leftmargin=1.5cm]
  \item Configurar um \textit{Identity Provider} (IdP) centralizado com Keycloak 26, incluindo \textit{realm} dedicado, clientes OIDC, papéis e grupos;
  \item Implementar autenticação OIDC/OAuth 2.0 com validação local de JWT (RS256/JWKS);
  \item Aplicar controlo de acesso baseado em papéis (RBAC) em endpoints de uma API REST;
  \item Activar e demonstrar MFA TOTP obrigatória para o papel de administrador;
  \item Implementar os três fluxos do ciclo de vida JML via Keycloak Admin REST API;
  \item Configurar auditoria de eventos de autenticação e administração;
  \item Containerizar toda a infraestrutura com Docker Compose para reprodutibilidade;
  \item Fornecer um dashboard interativo como evidência de demonstração.
\end{enumerate}

\subsection{Caso de Uso}

O cenário implementado simula uma \textbf{aplicação interna empresarial} com três perfis de utilizador:

\begin{itemize}[leftmargin=1.5cm]
  \item \textbf{Visitante}: acesso apenas a recursos públicos e ao próprio perfil autenticado;
  \item \textbf{Colaborador}: acesso a dados e funcionalidades internas de nível intermédio;
  \item \textbf{Administrador}: acesso total, incluindo gestão de utilizadores, auditoria e áreas MFA.
\end{itemize}

% ============================================================
%  2. ARQUITETURA
% ============================================================
\section{Arquitetura do Sistema}

\subsection{Visão Geral}

A arquitectura segue um modelo de \textbf{separação de responsabilidades} entre três componentes principais, todos orquestrados por Docker Compose:

\begin{figure}[H]
  \centering
  \includegraphics[width=0.92\textwidth]{arquitetura_iam.png}
  \caption{Arquitectura executável do projeto: FastAPI, Keycloak e PostgreSQL sob Docker Compose.}
  \label{fig:arquitetura}
\end{figure}

\begin{callout}[Fluxo principal]
  Cliente autentica no Keycloak $\rightarrow$ recebe JWT assinado (RS256) $\rightarrow$ apresenta token à FastAPI $\rightarrow$ FastAPI valida via JWKS localmente $\rightarrow$ aplica RBAC $\rightarrow$ responde.
\end{callout}

\subsection{Componentes e Tecnologias}

\begin{table}[H]
  \centering
  \caption{Componentes da arquitectura e respectivas funções}
  \begin{tabularx}{\textwidth}{@{}lXl@{}}
    \toprule
    \textbf{Componente} & \textbf{Função} & \textbf{Porto} \\
    \midrule
    Keycloak 26 & \textit{Identity Provider}: autenticação, emissão de JWT, realm, RBAC, MFA, JML, auditoria & 8081 \\
    \addlinespace
    FastAPI (Python 3.12) & Servidor de recursos: endpoints protegidos, validação JWT, RBAC, dashboard & 8000 \\
    \addlinespace
    PostgreSQL 16 & Base de dados persistente do Keycloak & 5432 \\
    \addlinespace
    JML Scripts (CLI) & Scripts Python para joiner/mover/leaver via Admin REST API & — \\
    \bottomrule
  \end{tabularx}
\end{table}

\subsection{Ficheiros de Configuração e Reprodutibilidade}

\begin{itemize}[leftmargin=1.5cm]
  \item \textbf{\texttt{docker-compose.yml}}: define os três serviços, redes internas, volumes e \textit{health checks};
  \item \textbf{\texttt{realm-export.json}}: exportação completa do realm \texttt{iam-tp} (papéis, grupos, utilizadores de teste, cliente OIDC, fluxos de autenticação MFA) — versionada em git;
  \item \textbf{\texttt{.env}}: variáveis de ambiente sensíveis (passwords, segredos do cliente OIDC) — não versionado.
\end{itemize}

\begin{lstlisting}[language=bash, caption=Arranque completo da infraestrutura]
# Clonar o repositório e criar o ficheiro .env
cp .env.example .env

# Iniciar todos os serviços (primeira vez: ~2-3 minutos)
docker compose up --build

# Verificar saúde dos serviços
docker compose ps
curl http://localhost:8000/health
# → {"status":"ok"}
\end{lstlisting}

\subsection{Estrutura de Ficheiros do Projeto}

\begin{verbatim}
Projeto_IAM/
├── docker-compose.yml          # Orquestração dos 3 serviços
├── realm-export.json           # Configuração versionável do realm
├── .env                        # Segredos (não versionado)
├── app/
│   ├── main.py                 # FastAPI: arranque e middlewares
│   ├── config.py               # Pydantic Settings
│   ├── auth.py                 # Validação JWT RS256, require_role()
│   ├── keycloak_client.py      # Admin API async helpers
│   └── routes/
│       ├── public.py           # /health, /me, /auth/token
│       ├── colaborador.py      # /colaborador/* (RBAC)
│       └── admin.py            # /admin/* (RBAC + MFA)
├── jml/
│   ├── joiner.py               # Criar utilizador
│   ├── mover.py                # Alterar papel + revogar sessões
│   └── leaver.py               # Desativar + revogar sessões
└── tests/
    └── test_auth.py            # Testes unitários (sem serviços)
\end{verbatim}

% ============================================================
%  3. AUTENTICAÇÃO
% ============================================================
\section{Autenticação OIDC/OAuth 2.0}

\subsection{Protocolo e Fluxo}

A autenticação é completamente delegada no Keycloak, seguindo o protocolo \textbf{OpenID Connect (OIDC)} sobre OAuth 2.0. A FastAPI actua como \textit{Resource Server}, validando os tokens emitidos pelo IdP.

Para a demonstração, é utilizado o fluxo \textit{Resource Owner Password Credentials} (ROPC) via o endpoint \texttt{/auth/token}. O realm está também configurado com \textit{Authorization Code + PKCE} para os fluxos browser recomendados em produção.

\begin{figure}[H]
  \centering
  \includegraphics[width=0.88\textwidth]{evidencias_dashboard.png}
  \caption{Dashboard do projeto: evidências de sessão, token JWT, claims OIDC e tabs de demonstração.}
  \label{fig:dashboard}
\end{figure}

\subsection{Validação JWT (RS256/JWKS)}

A validação dos tokens é feita \textbf{localmente na FastAPI}, sem chamadas ao Keycloak em cada pedido:

\begin{enumerate}[leftmargin=1.5cm]
  \item No primeiro pedido autenticado, a FastAPI descarrega a JWKS publicada pelo realm;
  \item A chave pública RS256 é cacheada em memória (\texttt{\_jwks\_cache} em \texttt{app/auth.py});
  \item Cada token é verificado localmente: assinatura, \texttt{issuer}, \texttt{audience} e expiração;
  \item As \textit{claims} relevantes são extraídas do payload para autorização.
\end{enumerate}

\begin{lstlisting}[language=python, caption=Lógica central de validação JWT em auth.py]
async def get_current_user(token: str = Depends(oauth2_scheme)) -> dict:
    jwks = await _get_jwks()          # cache em memória
    payload = jwt.decode(
        token,
        jwks,
        algorithms=["RS256"],
        audience=settings.oidc_client_id,
        options={"verify_exp": True},
    )
    return payload                    # claims disponíveis nos handlers
\end{lstlisting}

\subsection{Claims Utilizadas}

\begin{table}[H]
  \centering
  \caption{Claims JWT extraídas e utilizadas pela aplicação}
  \begin{tabularx}{\textwidth}{@{}lXl@{}}
    \toprule
    \textbf{Claim} & \textbf{Descrição} & \textbf{Uso} \\
    \midrule
    \texttt{sub} & Identificador único do utilizador (UUID Keycloak) & Auditoria, perfil \\
    \texttt{email} & Endereço de correio electrónico & Exibição no perfil \\
    \texttt{preferred\_username} & Nome de utilizador de login & Identificação visual \\
    \texttt{realm\_access.roles} & Lista de papéis do realm & Base do RBAC \\
    \texttt{amr} / \texttt{acr} & Método de autenticação (inclui MFA) & Verificação MFA \\
    \texttt{exp} & Timestamp de expiração & Validade do token \\
    \bottomrule
  \end{tabularx}
\end{table}

\subsection{Utilizadores de Teste}

\begin{table}[H]
  \centering
  \caption{Utilizadores de teste pré-configurados no realm \texttt{iam-tp}}
  \begin{tabular}{@{}llll@{}}
    \toprule
    \textbf{Username} & \textbf{Password} & \textbf{Papel} & \textbf{MFA} \\
    \midrule
    \texttt{admin.user} & \texttt{Admin@1234} & admin & TOTP obrigatório \\
    \texttt{colaborador.user} & \texttt{Colab@1234} & colaborador & Não configurado \\
    \texttt{visitante.user} & \texttt{Visit@1234} & visitante & Não configurado \\
    \bottomrule
  \end{tabular}
\end{table}

% ============================================================
%  4. AUTORIZAÇÃO RBAC
% ============================================================
\section{Autorização Baseada em Papéis (RBAC)}

\subsection{Modelo de Papéis}

O modelo RBAC implementado define três papéis de realm no Keycloak, com grupos correspondentes para facilitar a gestão:

\begin{itemize}[leftmargin=1.5cm]
  \item \textbf{\texttt{admin}} — Grupo \texttt{/Admins}: acesso total a todos os endpoints, incluindo administração e auditoria. MFA obrigatória.
  \item \textbf{\texttt{colaborador}} — Grupo \texttt{/Colaboradores}: acesso a dados e serviços internos.
  \item \textbf{\texttt{visitante}} — Grupo \texttt{/Visitantes}: acesso apenas a recursos públicos e perfil próprio.
\end{itemize}

\subsection{Matriz de Controlo de Acesso}

\begin{figure}[H]
  \centering
  \includegraphics[width=0.90\textwidth]{matriz_rbac.png}
  \caption{Matriz RBAC: mapeamento de endpoints para papéis com resultados esperados (200/401/403).}
  \label{fig:rbac}
\end{figure}

\begin{table}[H]
  \centering
  \caption{Controlo de acesso por endpoint}
  \rowcolors{2}{gray!10}{white}
  \begin{tabularx}{\textwidth}{@{}lXccc@{}}
    \toprule
    \textbf{Endpoint} & \textbf{Proteção} & \textbf{admin} & \textbf{colaborador} & \textbf{visitante} \\
    \midrule
    \texttt{GET /public} & Sem autenticação & \textcolor{green!60!black}{200} & \textcolor{green!60!black}{200} & \textcolor{green!60!black}{200} \\
    \texttt{GET /me} & JWT válido & \textcolor{green!60!black}{200} & \textcolor{green!60!black}{200} & \textcolor{green!60!black}{200} \\
    \texttt{GET /colaborador/data} & colaborador ou admin & \textcolor{green!60!black}{200} & \textcolor{green!60!black}{200} & \textcolor{red}{403} \\
    \texttt{GET /colaborador/perfil} & colaborador ou admin & \textcolor{green!60!black}{200} & \textcolor{green!60!black}{200} & \textcolor{red}{403} \\
    \texttt{GET /admin/users} & admin & \textcolor{green!60!black}{200} & \textcolor{red}{403} & \textcolor{red}{403} \\
    \texttt{GET /admin/audit} & admin & \textcolor{green!60!black}{200} & \textcolor{red}{403} & \textcolor{red}{403} \\
    \texttt{GET /admin/mfa-area} & admin + MFA & \textcolor{green!60!black}{200*} & \textcolor{red}{403} & \textcolor{red}{403} \\
    \bottomrule
    \multicolumn{5}{l}{\small * apenas se o token incluir prova de MFA (\texttt{amr=otp})}
  \end{tabularx}
\end{table}

\subsection{Implementação Técnica}

A autorização é aplicada através de dependências FastAPI reutilizáveis:

\begin{lstlisting}[language=python, caption=Dependências de autorização em auth.py]
def require_role(*roles: str):
    """Retorna 403 se o utilizador não tiver nenhum dos papéis indicados."""
    async def checker(user: dict = Depends(get_current_user)):
        user_roles = user.get("realm_access", {}).get("roles", [])
        if not any(r in user_roles for r in roles):
            raise HTTPException(status_code=403, detail="Acesso negado")
        return user
    return checker

def require_mfa():
    """Retorna 403 se o token não incluir prova de MFA."""
    async def checker(user: dict = Depends(get_current_user)):
        amr = user.get("amr", [])
        if not any(m in amr for m in ["otp", "mfa", "2fa"]):
            raise HTTPException(status_code=403, detail="MFA obrigatória")
        return user
    return checker
\end{lstlisting}

\begin{lstlisting}[language=python, caption=Aplicação das dependências nos endpoints]
# Apenas admin
@router.get("/admin/users")
async def list_users(user=Depends(require_role("admin"))):
    ...

# Admin com MFA comprovada
@router.get("/admin/mfa-area")
async def mfa_zone(user=Depends(require_mfa())):
    ...
\end{lstlisting}

% ============================================================
%  5. MFA
% ============================================================
\section{Autenticação Multifator (MFA/TOTP)}

\subsection{Configuração no Keycloak}

A MFA foi implementada com \textbf{TOTP (Time-based One-Time Password)}, compatível com aplicações como Google Authenticator, Microsoft Authenticator ou Authy.

A configuração inclui:

\begin{enumerate}[leftmargin=1.5cm]
  \item O fluxo de autenticação \textit{browser} do realm inclui a execução \texttt{auth-otp-form};
  \item O utilizador \texttt{admin.user} tem a \textit{required action} \texttt{CONFIGURE\_TOTP} activa;
  \item No primeiro acesso via browser, o Keycloak força a configuração do autenticador TOTP;
  \item Após configuração, o JWT inclui \texttt{"amr": ["otp"]} como prova de segundo factor.
\end{enumerate}

\subsection{Validação de MFA na API}

O endpoint \texttt{/admin/mfa-area} exige, cumulativamente:
\begin{itemize}[leftmargin=1.5cm]
  \item Papel \texttt{admin} em \texttt{realm\_access.roles};
  \item Presença de \texttt{otp}, \texttt{mfa} ou \texttt{2fa} no claim \texttt{amr} do JWT.
\end{itemize}

\begin{warningbox}[Evidência de demonstração]
  Token \texttt{admin} obtido via ROPC (sem MFA) $\rightarrow$ \texttt{/admin/mfa-area} retorna \textbf{403}.\\
  Token \texttt{admin} obtido via browser flow com TOTP $\rightarrow$ \texttt{/admin/mfa-area} retorna \textbf{200}.
\end{warningbox}

% ============================================================
%  6. JML
% ============================================================
\section{Ciclo de Vida de Identidades (JML)}

\subsection{Visão Geral}

O ciclo de vida JML (Joiner–Mover–Leaver) foi implementado sobre a \textbf{Keycloak Admin REST API}, garantindo que toda a gestão de identidades é centralizada no IdP e não em ficheiros locais ou bases de dados secundárias.

\begin{figure}[H]
  \centering
  \includegraphics[width=0.90\textwidth]{fluxo_jml.png}
  \caption{Fluxos JML: criação, alteração de papel/grupo, desactivação e revogação de sessões.}
  \label{fig:jml}
\end{figure}

\subsection{Joiner — Integração de Novo Utilizador}

\begin{itemize}[leftmargin=1.5cm]
  \item Cria o utilizador no realm com atributos, credenciais e estado activo;
  \item Atribui o papel de realm inicial (\textit{visitante}, \textit{colaborador} ou \textit{admin});
  \item Adiciona o utilizador ao grupo correspondente;
  \item Disponível como CLI (\texttt{jml/joiner.py}) e como endpoint REST (\texttt{POST /admin/jml/joiner}).
\end{itemize}

\begin{lstlisting}[language=bash, caption=Exemplo de criação via CLI]
python jml/joiner.py \
  --username novo.colaborador \
  --email novo@empresa.pt \
  --role colaborador \
  --password "Senha@2026"
\end{lstlisting}

\subsection{Mover — Alteração de Papel}

\begin{itemize}[leftmargin=1.5cm]
  \item Remove o papel e grupo anterior do utilizador;
  \item Atribui o novo papel e grupo;
  \item \textbf{Revoga todas as sessões activas} via \texttt{DELETE /users/\{id\}/sessions}, forçando reautenticação imediata com os novos privilégios.
\end{itemize}

\subsection{Leaver — Saída do Utilizador}

\begin{itemize}[leftmargin=1.5cm]
  \item Desactiva a conta (\texttt{enabled: false}) — impede novos logins;
  \item Remove todos os papéis de realm;
  \item \textbf{Revoga todas as sessões activas} — corta acesso imediatamente, sem aguardar expiração do token.
\end{itemize}

\subsection{Estratégia de Revogação de Sessões}

\begin{callout}[Porquê a revogação explícita?]
  Sem revogação explícita, um utilizador movido ou desactivado manteria o acesso durante o tempo de vida do JWT (por omissão, 300 segundos). A chamada \texttt{DELETE /users/\{id\}/sessions} força a reautenticação imediata, garantindo que a alteração de papel ou desactivação tem efeito instantâneo.
\end{callout}

\subsection{Tabela Resumo JML}

\begin{table}[H]
  \centering
  \caption{Resumo das operações JML e efeitos de segurança}
  \begin{tabularx}{\textwidth}{@{}lXl@{}}
    \toprule
    \textbf{Fluxo} & \textbf{Acções realizadas} & \textbf{Efeito imediato} \\
    \midrule
    Joiner & Cria utilizador, atribui papel e grupo & Acesso inicial controlado \\
    Mover & Troca papel/grupo, revoga sessões & Novo acesso após reautenticação \\
    Leaver & Desactiva conta, remove papéis, revoga sessões & Bloqueio imediato \\
    \bottomrule
  \end{tabularx}
\end{table}

% ============================================================
%  7. AUDITORIA E SEGURANÇA
% ============================================================
\section{Auditoria e Segurança Operacional}

\subsection{Eventos de Auditoria}

O Keycloak regista dois tipos de eventos:

\begin{enumerate}[leftmargin=1.5cm]
  \item \textbf{Eventos de utilizador}: login, falha de login, logout, refresh de token, configuração e falha de OTP;
  \item \textbf{Eventos administrativos}: criação/remoção de utilizadores, alteração de papéis, actualização de grupos.
\end{enumerate}

A API expõe o endpoint \texttt{GET /admin/audit} (protegido por papel \texttt{admin}) que consulta os eventos via Admin REST API. O dashboard apresenta os eventos com filtragem por tipo e um resumo visual por categoria.

\subsection{Análise de Riscos e Mitigações}

\begin{table}[H]
  \centering
  \caption{Mapeamento de riscos e mitigações implementadas}
  \rowcolors{2}{gray!10}{white}
  \begin{tabularx}{\textwidth}{@{}lX@{}}
    \toprule
    \textbf{Risco} & \textbf{Mitigação Implementada} \\
    \midrule
    Token inválido ou forjado & Validação RS256 via JWKS com verificação de \texttt{issuer} e \texttt{audience} \\
    Acesso indevido a recursos & RBAC no backend com \texttt{require\_role()} por endpoint \\
    Acesso administrativo sem segundo factor & TOTP obrigatório (\textit{required action}) e validação de \texttt{amr} na API \\
    Sessões activas após mudança de papel & Revogação explícita de sessões no Mover e Leaver \\
    Credenciais em ficheiros de configuração & Variáveis de ambiente via \texttt{.env} (não versionado); secrets em Docker Compose \\
    Falta de rastreabilidade & Eventos Keycloak e endpoint de auditoria protegido \\
    \bottomrule
  \end{tabularx}
\end{table}

\subsection{Configuração de Segurança do Realm}

\begin{itemize}[leftmargin=1.5cm]
  \item \textbf{Algoritmo de assinatura}: RS256 com par de chaves gerido pelo Keycloak;
  \item \textbf{Tempo de vida do token}: 300 segundos (access token), 30 minutos (refresh token);
  \item \textbf{PKCE}: activado para o cliente \texttt{fastapi-client} (fluxo browser);
  \item \textbf{Direct Access Grants}: activado apenas para testes (ROPC);
  \item \textbf{Redirect URIs}: restritas a \texttt{localhost:8000} e \texttt{localhost:8081};
  \item \textbf{Eventos de auditoria}: activados para todos os tipos de eventos de utilizador e administração.
\end{itemize}

% ============================================================
%  8. TESTES E VALIDAÇÃO
% ============================================================
\section{Testes e Validação}

\subsection{Testes Unitários}

Foram implementados quatro testes unitários em \texttt{tests/test\_auth.py} que validam a lógica central de autenticação e autorização sem necessitar de serviços em execução:

\begin{table}[H]
  \centering
  \caption{Cobertura dos testes unitários}
  \begin{tabularx}{\textwidth}{@{}lX@{}}
    \toprule
    \textbf{Teste} & \textbf{O que valida} \\
    \midrule
    \texttt{test\_require\_role\_allows\_admin} & \texttt{require\_role("admin")} aceita token com papel admin \\
    \texttt{test\_require\_role\_denies\_non\_admin} & Papel incorrecto retorna HTTP 403 \\
    \texttt{test\_require\_mfa\_denies\_without\_otp} & Ausência do claim \texttt{amr=otp} retorna HTTP 403 \\
    \texttt{test\_get\_current\_user\_decodes\_rs256} & Validação de assinatura RS256, issuer e audience com JWKS simulado \\
    \bottomrule
  \end{tabularx}
\end{table}

Os testes geram um par RSA temporário para simular as chaves do Keycloak, usando \textit{mocking} do método \texttt{\_get\_jwks()}.

\begin{lstlisting}[language=bash, caption=Execução dos testes unitários]
pip install pytest
pytest tests/ -v
\end{lstlisting}

\subsection{Validação Funcional via Dashboard}

O dashboard interativo (\texttt{http://localhost:8000/dashboard}) permite validar visualmente todos os requisitos:

\begin{enumerate}[leftmargin=1.5cm]
  \item \textbf{Aba Autenticação}: formulário de login, descodificador de JWT, visualização de claims;
  \item \textbf{Aba Autorização}: matriz RBAC com botão ``Testar Todos os Endpoints'' — demonstra 200/403 automaticamente;
  \item \textbf{Aba JML}: formulários Joiner, Mover e Leaver com chamadas à API em tempo real;
  \item \textbf{Aba Auditoria}: lista de eventos filtráveis por tipo, resumo estatístico;
  \item \textbf{Aba MFA}: demonstração do fluxo TOTP e do endpoint \texttt{/admin/mfa-area};
  \item \textbf{Aba Arquitectura}: diagrama de componentes e mapeamento de portos.
\end{enumerate}

% ============================================================
%  9. CONCLUSÃO
% ============================================================
\section{Conclusão}

\subsection{Resultados Obtidos}

O projeto implementa com sucesso uma plataforma IAM completa que cobre todos os requisitos do trabalho prático:

\begin{itemize}[leftmargin=1.5cm]
  \item[$\checkmark$] Keycloak como IdP centralizado com realm dedicado (\texttt{iam-tp});
  \item[$\checkmark$] Autenticação OIDC/OAuth 2.0 com validação JWT RS256/JWKS;
  \item[$\checkmark$] RBAC com três papéis e grupos sincronizados;
  \item[$\checkmark$] MFA TOTP obrigatória para administradores, validada por \texttt{amr};
  \item[$\checkmark$] JML completo (Joiner, Mover, Leaver) com revogação imediata de sessões;
  \item[$\checkmark$] Auditoria de eventos de autenticação e administração;
  \item[$\checkmark$] Infraestrutura containerizada e reproduzível com \texttt{docker compose up};
  \item[$\checkmark$] Dashboard interativo como prova de demonstração.
\end{itemize}

\subsection{Decisões Técnicas Relevantes}

\begin{itemize}[leftmargin=1.5cm]
  \item A \textbf{validação local de JWT} (cache JWKS) elimina a dependência do Keycloak em cada pedido autenticado, reduzindo latência e aumentando resiliência;
  \item A \textbf{revogação explícita de sessões} no Mover e Leaver garante que alterações de acesso têm efeito imediato, sem aguardar a expiração natural do token;
  \item A \textbf{exportação do realm} em \texttt{realm-export.json} permite reproduzir toda a configuração IAM com um único comando, tornando o ambiente completamente declarativo.
\end{itemize}

\subsection{Aprendizagens}

Este projecto permitiu consolidar os conceitos fundamentais de IAM na prática: a diferença entre autenticação e autorização, o papel central do IdP num ecossistema federado, as implicações de segurança da gestão de sessões e tokens, e os desafios operacionais do ciclo de vida de identidades em ambientes reais.

% ============================================================
%  REFERÊNCIAS
% ============================================================
\newpage
\section*{Referências}
\addcontentsline{toc}{section}{Referências}

\begin{enumerate}[leftmargin=1.5cm]

  \item Keycloak Project. \textit{Keycloak 26 Documentation}. Red Hat, 2024. Disponível em: \url{https://www.keycloak.org/documentation}

  \item Hardt, D. \textit{The OAuth 2.0 Authorization Framework}. RFC 6749, IETF, 2012. Disponível em: \url{https://tools.ietf.org/html/rfc6749}

  \item Sakimura, N. et al. \textit{OpenID Connect Core 1.0}. OpenID Foundation, 2014. Disponível em: \url{https://openid.net/specs/openid-connect-core-1_0.html}

  \item Jones, M. et al. \textit{JSON Web Token (JWT)}. RFC 7519, IETF, 2015. Disponível em: \url{https://tools.ietf.org/html/rfc7519}

  \item Jones, M. et al. \textit{JSON Web Key (JWK)}. RFC 7517, IETF, 2015. Disponível em: \url{https://tools.ietf.org/html/rfc7517}

  \item FastAPI. \textit{FastAPI Documentation}. Sebastián Ramírez, 2024. Disponível em: \url{https://fastapi.tiangolo.com/}

  \item Docker Inc. \textit{Docker Compose Documentation}. 2024. Disponível em: \url{https://docs.docker.com/compose/}

  \item M'Raihi, D. et al. \textit{TOTP: Time-Based One-Time Password Algorithm}. RFC 6238, IETF, 2011. Disponível em: \url{https://tools.ietf.org/html/rfc6238}

  \item NIST. \textit{Digital Identity Guidelines — Authentication and Lifecycle Management}. SP 800-63B, 2017.

  \item Costa, R. \textit{Gestão de Identidade — Materiais das aulas}. ESTG-IPP, MEI, 2025/2026.

\end{enumerate}

\end{document}
